\section*{Abstract}
Accurate identification of skin conditions is important for understanding skin health and supporting appropriate care. Deep learning techniques can classify skin conditions from images, while transfer learning enables the use of pre-trained models for this task. However, model performance may vary depending on the network architecture and the division of data into training, validation, and testing sets. Furthermore, visual explanations can help identify image regions that influence model predictions. This study evaluates EfficientNetB0, MobileNetV2, and ResNet50 for image-based skin condition classification and severity identification, using Gradient-weighted Class Activation Mapping (Grad-CAM) to visualize regions associated with their predictions. The developed system classifies images into 12 labels across four categories: Acne, Burn, Eczema, and Normal skin. The first three categories include mild, moderate, and severe levels, whereas the Normal category comprises dry, normal, and oily skin types. The models are evaluated using different dataset splitting ratios and compared based on classification performance and inference time. Among the evaluated configurations, ResNet50 with a 90\% training, 5\% validation, and 5\% testing split achieved the highest recorded accuracy of 86.84\%, precision of 87.22\%, recall of 86.84\%, F1-score of 86.79\%, and ROC-AUC of 99.27\%, with an inference time of 72.42 ms per image. Grad-CAM provides visual explanations by highlighting image regions associated with model predictions. These findings demonstrate the potential of transfer learning and explainable deep learning for skin condition classification and severity assessment, while highlighting the need for evaluation on more diverse and independent datasets to assess generalization to unseen images.
\\
\vspace{1cm}

\textbf{Keywords:} Skin condition classification; Transfer learning; Deep learning; Explainable artificial intelligence (XAI); Grad-CAM; Convolutional neural networks
\pagebreak
